\chapter{Dimensionnement}
\label{chap:dimensionnement}

\begin{todobox}[Chapitre à compléter après audit]
Les modèles exacts de CPU, le nombre de disques et la RAM installée des R630, ainsi que l'inventaire des serveurs historiques, sont nécessaires pour finaliser ce chapitre.
\end{todobox}

\section{Méthodologie}
Le dimensionnement suit une logique N+1 : la perte d'un nœud ne doit pas entraîner la perte définitive des workloads. Le dimensionnement utile est calculé après réservation des ressources système et après tolérance à la panne.

\section{Formules}
\begin{itemize}
  \item \textbf{Capacité brute} : somme des ressources de tous les nœuds.
  \item \textbf{Capacité système} : réservation hyperviseur (10–15\% CPU, 4–8 Go RAM par nœud).
  \item \textbf{Capacité utile N+1} : capacité système $\times$ (N-1)/N.
  \item \textbf{Overcommit} : ratio CPU (2–4$\times$) et RAM (1,0–1,5$\times$).
\end{itemize}

\section{Hypothèses de profils VM}
\begin{tabularx}{\textwidth}{@{}lrrrX@{}}
\toprule
\textbf{Profil} & \textbf{vCPU} & \textbf{RAM} & \textbf{Disque} & \textbf{Exemple} \\
\midrule
Petite Linux      & 1 & 2 Go  & 20 Go  & TP, dev \\
Standard Linux    & 2 & 4 Go  & 40 Go  & App web, DB légère \\
Développement     & 4 & 8 Go  & 60 Go  & IDE, build \\
Base de données   & 4 & 16 Go & 100 Go & PostgreSQL, MariaDB \\
Kubernetes / CI   & 4–8 & 16 Go & 100 Go & Cluster de test \\
\bottomrule
\end{tabularx}

\section{Scénarios à chiffrer}
\begin{todobox}[À produire]
\begin{itemize}
  \item Scénario minimal : 4 R630, X Go RAM, Y disques, Z VMs.
  \item Scénario cible : serveurs historiques, capacité annoncée 500 utilisateurs.
  \item Scénario intermédiaire : ajout de RAM et disques aux R630.
\end{itemize}
\end{todobox}